A name's meaning in \PurePy must not depend on the order in which modules load or on runtime state.
\figref{load-dependent} lists the import and module behaviours of Python excluded for this reason, and says
how and when \PurePy excludes each.

\begin{figure}
\begin{AddedBlock}
\centering
\footnotesize
\setlength{\tabcolsep}{4pt}
\begin{tabular}{>{\raggedright\arraybackslash}p{0.24\textwidth}>{\raggedright\arraybackslash}p{0.27\textwidth}>{\raggedright\arraybackslash}p{0.10\textwidth}>{\raggedright\arraybackslash}p{0.30\textwidth}}
\arrayrulecolor{gray!60}\hline
\rowcolor{grammarheadcolor}Construct & Python & Exclusion & Mechanism \\
Import inside function body or branch, or after a statement & Dependencies of module known only as its statements run & Syntactic & Module body is an import prefix followed by a top-level statement \\\rowgap
Wildcard import \pyinline{from q import *} & Introduces names of \pyinline{q} implicitly & Syntactic & Not in grammar \\\rowgap
From-import of submodule through its parent, \pyinline{from a import b} in module \pyinline{a.b} & Parent may still be loading & Static & From-imported submodule is a dependency, so the import is a cycle \\\rowgap
Import of contained module, \pyinline{import x.y} in module \pyinline{x} & Binding for head denotes still-loading module & Static & Containing modules of a plain import are checked, importing module included \\\rowgap
Reference \pyinline{pkg.sub.y} after \pyinline{import pkg} alone & Resolves if any module has loaded \pyinline{pkg.sub}; \pyinline{AttributeError} otherwise & Static & Qualified names resolve through imports of referring module \\\rowgap
Reference to module member assigned on some paths only & \pyinline{AttributeError} if path not taken & Static & Member must be definitely assigned \\\rowgap
Module binding name of one of its submodules & Attribute set by assignment and by loading submodule; last wins & Static & Names bound by module disjoint from names of its submodules
\end{tabular}
\caption{Python import and module behaviours excluded by \PurePy}
\label{fig:load-dependent}
\end{AddedBlock}
\end{figure}


The last three exclusions rule out programs whose meaning under Python changes when a module elsewhere in
the program is loaded. A submodule is not an attribute of its parent until it is loaded, so after
\pyinline{import pkg} alone, Python resolves \pyinline{pkg.sub.y} if some earlier import, anywhere in the
program, has loaded \pyinline{pkg.sub}, and raises \pyinline{AttributeError} otherwise. An attribute of a
module can be set both by an assignment in the module body and by loading a same-named submodule, each
overwriting the other, so the value of the attribute depends on load order. Consider a program with four
modules:

\begin{minipage}[t]{0.2\textwidth}
\begin{python}
# pkg
sub: int = 1
\end{python}
\end{minipage}
\hfill
\begin{minipage}[t]{0.2\textwidth}
\begin{python}
# pkg.sub
# (empty)
\end{python}
\end{minipage}
\hfill
\begin{minipage}[t]{0.25\textwidth}
\begin{python}
# helper
import pkg.sub
\end{python}
\end{minipage}
\hfill
\begin{minipage}[t]{0.28\textwidth}
\begin{python}
# __main__
import pkg
print(pkg.sub) # 1
\end{python}
\end{minipage}

Here \pyinline{helper} is never imported, so \pyinline{pkg.sub} never loads, and \pyinline{pkg.sub} in the
main module denotes the assignment. Now let the main module also import \pyinline{helper}, which it does not
otherwise use:

\begin{python}
# __main__
import helper
import pkg
print(pkg.sub) # <module 'pkg.sub'>
\end{python}

Loading \pyinline{pkg.sub}, wherever in the program it happens, sets it as an attribute of the shared
\pyinline{pkg} object, so the extra import changes the meaning of \pyinline{pkg.sub} in the main module, even though
it does not import \pyinline{pkg.sub}. \PurePy rejects \pyinline{pkg}, because it binds the name of its submodule.
